\appendix
\section{Three-row and graph-triple estimates}\label{app:section}
This appendix proves the triple estimate $H_T^2\succeq H_T$ used in
Proposition~\ref{asm:global}. The three-row variance inequality,
Lemma~\ref{app:rows}, is the conditional-projection inequality of Fu, Qin and
Wang~\cite[Theorem~5.1]{FQW26}, in the equivalent form \eqref{app:rowvariance}.
We reproduce the alternative proof of OpenAI~\cite[Section~3]{OAI26}. Its
fixed-block step, Lemma~\ref{app:fixed}, uses transposition eigenspaces, a
strategy that~\cite[Section~3]{OAI26} credits to the symmetric-group method
of Carlen, Lieb and Loss~\cite[Section~2]{CLL06}. Its count step,
Lemma~\ref{app:counts}, uses an eigenvector increment argument with the
adjacent-count coupling of Fu, Qin and Wang~\cite[Lemma~5.2]{FQW26}; the count
law corresponds to their Lemmas~5.2--5.4. The transfer to graph triples,
Lemma~\ref{app:triple}, uses the triple encoding of
OpenAI~\cite[Proposition~4.1]{OAI26}. All steps are included.

\subsection{A variance formulation}
For three coordinates on any finite probability space let $V_i$ be the
centered functions of coordinate $i$, and put the sum of marginal squared
norms on $V_1\oplus V_2\oplus V_3$. For
$A(g_1,g_2,g_3)=g_1+g_2+g_3$ the adjoint on centered joint functions is
$A^*f=(P_1f,P_2f,P_3f)$, where $P_i$ is conditional expectation onto that
coordinate. This follows by testing inner products. Therefore the assertions
\begin{equation}\label{app:rowvariance}
 \operatorname{Var}\Big(\sum_{i=1}^3g_i\Big)\le2\sum_{i=1}^3\operatorname{Var}(g_i),
 \qquad \sum_{i=1}^3P_i\preceq2I\text{ on }\mathbf 1^\perp
\end{equation}
are equivalent, since $AA^*=\sum_iP_i$ and both assertions say
$\lVert A\rVert^2\le2$.

\subsection{Fixed block sizes}
\begin{lemma}[Fixed block sizes]\label{app:fixed}
Take independent uniform ordered partitions of labeled sets $U,W$ into three
blocks of prescribed sizes, and let the row coordinates be $(U_i,W_i)$. Then
\eqref{app:rowvariance} holds for these coordinates.
\end{lemma}
\begin{proof}
Let $J$ be the sum of $I-T_\tau$ over label transpositions within either set.
These are self-adjoint involutions, and
$\langle f,Jf\rangle=\tfrac12\sum_\tau\lVert f-T_\tau f\rVert^2$.
The two permutation groups act transitively on the partition pairs, so
$\ker J=\mathbb R\mathbf 1$, including the degenerate one-state case.

Every transposition preserves each row-function space: it takes $U_i$ or
$W_i$ to its image under the known label permutation, without needing the
other blocks. Thus $J$ preserves these spaces. In finite dimension its
spectral projectors are polynomials in $J$ (interpolate on the distinct
eigenvalues), so they also preserve the row spaces. Center each $g_i$ and
project it into each positive eigenspace, writing $g_{i,\lambda}$.
At a given state any transposition changes at most two block tuples; hence
at most two of the three differences $g_{i,\lambda}-T_\tau g_{i,\lambda}$
are nonzero. The inequality $(a+b)^2\le2(a^2+b^2)$ gives
\[
 \lambda\Big\lVert\sum_i g_{i,\lambda}\Big\rVert^2
 \le2\lambda\sum_i\lVert g_{i,\lambda}\rVert^2.
\]
Cancel $\lambda>0$ and sum over orthogonal eigenspaces to obtain
\eqref{app:rowvariance}. The invariance of the row spaces is what permits
passing from the transposition energy estimate to this variance estimate.
\end{proof}

\subsection{Random block sizes}
\begin{lemma}[Random block sizes]\label{app:counts}
For nonnegative integers $m$ and $\alpha_1,\alpha_2,\alpha_3$ let
$w_i(t)=1/(t!(t+\alpha_i)!)$, and condition the product weights
$\prod_iw_i(x_i)$ on $x_1+x_2+x_3=m$. Then \eqref{app:rowvariance} holds for
the coordinates $x_i$.
\end{lemma}
\begin{proof}
All compositions of $m$ have positive probability. First consider two
coordinates of total $s$, with $p_s(y)\propto w_j(y)w_k(s-y)$,
$0\le y\le s$. On $0\le y<s$,
\[
 \frac{p_s(y)}{p_{s-1}(y)}
 =\frac{c_s}{(s-y)(s-y+\alpha_k)},\qquad c_s>0,
\]
which increases with $y$. Reweighting by an increasing positive likelihood
ratio increases expectations of increasing functions: for increasing $a,b$
and independent $Y,Y'$ with the old law,
\[
 2\operatorname{Cov}(a(Y),b(Y))
 =\mathbb E[(a(Y)-a(Y'))(b(Y)-b(Y'))]\ge0.
\]
The extra mass at the new largest value $s$ preserves this stochastic order,
so $Y_{s-1}\le_{\rm st}Y_s$. Apply the same argument to $Z_s=s-Y_s$ to obtain
$Y_s\le_{\rm st}Y_{s-1}+1$. Common-uniform quantiles simultaneously couple
these variables as $Y_{s-1}\le Y_s\le Y_{s-1}+1$. Thus increasing the total
by one increases exactly one of the two integer coordinates by one.

Use the addition operator $A$ above for the centered marginal spaces.
A nonzero eigenvector $b$ of $A^*A$ at eigenvalue $\lambda$ satisfies
\[
 \lambda b_i(t)=b_i(t)+\sum_{j\ne i}\mathbb E[b_j(x_j)\mid x_i=t],
 \qquad 0\le t\le m.
\]
If $m>0$, set $\eta=\max_{i,t<m}|b_i(t+1)-b_i(t)|>0$; otherwise centering
would make $b$ zero. Adjacent values of $x_i$ change the other total by one.
The coupling just proved changes exactly one of those coordinates by one,
so the difference of their summed conditional expectations has absolute
value at most $\eta$. Subtracting the eigenvector equations gives
\[
 |\lambda-1|\,|b_i(t+1)-b_i(t)|\le\eta.
\]
At a maximizing increment, $|\lambda-1|\le1$ and hence $\lambda\le2$.
Apply this to a largest eigenvalue of the positive $A^*A$ to prove
\eqref{app:rowvariance}. If $m=0$, every variance is zero.
\end{proof}

\subsection{All binary three-row margins}
\begin{lemma}[Three-row variance inequality]\label{app:rows}
Take the uniform binary matrices with any feasible three-row margins. Then
\eqref{app:rowvariance} holds for the three row coordinates.
\end{lemma}
\begin{proof}
Delete columns of sum zero or three and subtract the sum-three count from
the row totals. Let $U$ be the sum-one columns and $W$ the sum-two columns.
Write $U_i$ for the unique-one locations in row $i$, $W_i$ for the
unique-zero locations, and $r_i$ for the adjusted row total. Each row is
equivalent to $(U_i,W_i)$, and the whole matrix to the two ordered
partitions. With $k_i=|W_i|$ and $\delta_i=r_i-|W|$, one has
$|U_i|=k_i+\delta_i$. At fixed $k$ there are exactly
\[
 \frac{|U|!\,|W|!}{\prod_i k_i!(k_i+\delta_i)!}
\]
matrices, and the two uniform partition laws are independent. All these
partitions are feasible by their unique-one and unique-zero interpretation.

Set $l_i=\max(0,-\delta_i)$, $x_i=k_i-l_i$, $\alpha_i=|\delta_i|$, and
$m=|W|-\sum_i l_i$. Nonemptiness gives $m\ge0$. The feasible $x$ are all
nonnegative compositions of $m$; block upper bounds follow from their
nonnegative sizes and fixed totals. Also
$k_i!(k_i+\delta_i)!=x_i!(x_i+\alpha_i)!$, so the count law is the one of
Lemma~\ref{app:counts}. For a row function $f_i$, the conditional mean
$b_i(k_i)=\mathbb E[f_i\mid k]$ depends only on $k_i$, because its marginal
blocks are independent uniform subsets of sizes $k_i+\delta_i,k_i$.
Total variance, Lemma~\ref{app:fixed}, and then Lemma~\ref{app:counts} give
\begin{align*}
 \operatorname{Var}\Big(\sum_i f_i\Big)
 &=\mathbb E\operatorname{Var}\Big(\sum_i f_i\Bigm| k\Big)
  +\operatorname{Var}\Big(\sum_i b_i\Big)\\
 &\le2\sum_i\mathbb E\operatorname{Var}(f_i\mid k)+2\sum_i\operatorname{Var}(b_i)
 =2\sum_i\operatorname{Var}(f_i).
\end{align*}
This proves \eqref{app:rowvariance} for every feasible three-row margin.
\end{proof}

\subsection{Encoding graph triples}
Fix $T=\{i,j,k\}$, the graph outside $T$, and each outside vertex's number of
neighbors in $T$. These data define cells preserved by pair updates inside
$T$. The degree sum in $T$ minus the sum of outside counts determines twice
the internal edge count $e$. Encode the outside incidences by three rows,
and augment them as in Table~\ref{app:encoding}, where $s_i$ is the internal
degree of vertex $i$.

\begin{table}[ht]
\centering
\begin{tabular}{c l l l}
\toprule
$e$ & Internal graph & Extra column & Prescribed row sum\\
\midrule
0 & empty & none & $d_i$\\
1 & edge between marked vertices & $s_i$, sum 2 & $d_i$\\
2 & path centered at marked vertex & $s_i-1$, sum 1 & $d_i-1$\\
3 & triangle & none & $d_i-2$\\
\bottomrule
\end{tabular}
\caption{Augmentation of the three-row encoding of a graph triple by its
internal edge count $e$.}
\label{app:encoding}
\end{table}

The extra column is a separate label and uniquely reconstructs the internal
graph. Conversely every binary matrix of these margins reconstructs a graph
in the cell: the row sums impose the three degrees, and the column sums and
fixed outside graph impose all other degrees. There is no extra symmetry or
structural-zero restriction. This is a bijection preserving uniform measure.

An augmented row $S_i$ fixes the external neighbors of $i$, hence its
internal degree $s_i$ by subtraction from $d_i$; its extra entry adds no
restriction. Together with the cell data it determines the full pair-fiber
signature for $\{j,k\}$: all edges away from that pair, its internal edge
$jk$ with indicator $e-s_i$, the other outside pair counts obtained by
subtracting incidences to $i$, and the pair count $s_i$ of vertex $i$ itself.
Conversely the signature determines the augmented row. In particular $s_i=1$
fixes a count, not a choice of $j$ or $k$. Both conditionings average exactly
the same sets, so $E_{\{j,k\}}=P_i=\mathbb E[\,\cdot\mid S_i]$ on the cell.

\begin{lemma}[Graph triples]\label{app:triple}
For every vertex triple $T$, $H_T^2\succeq H_T$.
\end{lemma}
\begin{proof}
Lemma~\ref{app:rows} and the adjoint formulation \eqref{app:rowvariance}
give $\sum_iP_i\preceq2I$ on centered cell functions. Thus
$H_T=3I-\sum_iP_i\succeq I$ there, while $H_T$ kills cell constants.
Self-adjointness makes its eigenvalues zero or at least one, and hence
$H_T^2\succeq H_T$. Taking the orthogonal direct sum over cells proves the
claim on the entire graph space, including one-state cells.
\end{proof}
